% Reconstructed text from V5 PDF page 29; non-native source
\par\smallskip\noindent\rule{\textwidth}{0.4pt}\par{\sffamily\footnotesize Original V5 page 29 -- recovered text}\par\begin{flushleft}\footnotesize\ttfamily
\par\smallskip
\par\smallskip
\noindent Diagnostics and Serving Infrastructure for MAS 29\par
\par\smallskip
\par\smallskip
\par\smallskip
\noindent SGLang.\par
\noindent  Leyline introduces declarative directive splicing for KV cache management: agents can\par
\noindent specify retention directives as a 4-tuple declaring edits and mode (in-place splice or prefix-\par
\noindent trimmed re-prefill for “semantic forgetting”). The architecture-agnostic interface routes each di-\par
\noindent rective to a per-architecture kernel with closed-form RoPE-rotation correction, bridging external\par
\noindent policy-level decisions with internal serving-system execution. Results show +11.2 percentage-\par
\noindent point replay cache-hit rate, up to 241 ms latency reduction, and +14.3 percentage-point agentic\par
\noindent solve rate on debug-gym. Leyline fills a specific gap where “kernel-level eviction methods make\par
\noindent their own decisions and cannot accept policy directives from outside the kernel” — enabling\par
\noindent agents to express retention preferences that current serving systems cannot accommodate.\par
\noindent  PolyKV implements a shared asymmetrically-compressed KV cache pool for multi-agent\par
\noindent inference. A single shared pool uses write-once, read-many semantics: keys are quantized\par
\noindent to int8, values are compressed to 3-bit via TurboQuant MSE (FWHT rotation + Lloyd-Max\par
\noindent quantization). Multiple HuggingFace DynamicCache objects read from the shared pool simul-\par
\noindent taneously. This achieves a stable 2.91× compression ratio with 97.7\% KV memory reduction\par
\noindent for 15 concurrent agents at 4K context, at only +0.57\% perplexity degradation and BERTScore\par
\noindent F1 of 0.928. Limitations include evaluation on only 2 model architectures, maximum 15 agents\par
\noindent tested, and a fixed lossy compression ratio.\par
\noindent  ConServe replaces prediction with observation for agentic serving scheduling. Rather than\par
\noindent estimating future tool-call durations (Continuum) or future prefix reuse (KVFlow), ConServe\par
\noindent uses only two directly observable metrics — per-decoder KV occupancy and first-turn input\par
\noindent length — to make conversation-level scheduling decisions. KV cache is transferred to the de-\par
\noindent code node once and pinned for the entire remaining conversation tail. This observation-based\par
\noindent approach achieves 51.08\% p95 TTFET reduction and +7.51\% energy improvement on homoge-\par
\noindent neous GPUs (+22.75\% on heterogeneous hardware). The design philosophy — that observation\par
\noindent is more robust than prediction for agentic workloads with inherently unpredictable tool-call pat-\par
\noindent terns — offers a counterpoint to the prediction-based approaches of KVFlow and Continuum.\par
\noindent  KAIROS [44] introduces context-growth-aware power management for multi-agent serving.\par
\noindent The system monitors per-agent context growth rate and GPU memory pressure, adjusting GPU\par
\noindent frequency and concurrency to balance power consumption against SLO compliance. When\par
\noindent memory headroom is sufficient, agents receive full-frequency service; when context growth\par
\noindent threatens to push the system into a cache-thrashing regime, KAIROS lowers GPU frequency\par
\noindent to reduce power without sacrificing throughput (since the workload becomes memory-bound).\par
\noindent This achieves an average 27\% power reduction (up to 39.8\%) while meeting latency SLOs.\par
\noindent KAIROS is the only infrastructure system that explicitly manages the power dimension, and its\par
\noindent control loop implicitly bridges monitoring state (context size, agent progress, memory pressure)\par
\noindent with serving decisions (frequency, concurrency, placement).\par
\noindent  LRAgent decomposes KV cache into a shared base component (computed from pretrained\par
\noindent weights) and adapter-dependent low-rank residuals, enabling efficient cache sharing across\par
\noindent multi-LoRA agents. A custom Flash-LoRA-Attention kernel reconstructs full attention outputs\par
\noindent without materializing the full-dimensional adapter-specific KV cache. This achieves throughput\par
\noindent and TTFT approaching the fully-shared caching upper bound while maintaining accuracy near\par
\noindent the non-shared (independent per-adapter) baseline. The approach requires a “shared-A” multi-\par
\noindent LoRA architecture where all adapters share the same LoRA-A matrix, and its custom kernel\par
\noindent limits portability to new architectures (ICML 2026).\par
\end{flushleft}
